%%% FIELD proof-causal-identification
Fix \(a\in\{0,1\}\), and define
\[
\rho_a(x):=P(A=a\mid X=x)
=
\begin{cases}
\pi(x),&a=1,\\
1-\pi(x),&a=0.
\end{cases}
\]
By \(\mathrm{ass{:}overlap}\), \(\rho_a(x)\ge e_0>0\) for every \(x\in[0,1]^d\). Because \(Y_a\) is bounded, let \(g_a(X)\) be a version of \(\mathbb E_P[Y_a\mid X]\). For every bounded Borel function \(b:[0,1]^d\to\mathbb R\), consistency gives
\[
\mathbb E_P\!\left[b(X)\mathbf 1\{A=a\}Y\right]
=\mathbb E_P\!\left[b(X)\mathbf 1\{A=a\}Y_a\right].
\]
Armwise conditional exchangeability then yields
\[
\begin{aligned}
\mathbb E_P\!\left[b(X)\mathbf 1\{A=a\}Y_a\right]
&=\mathbb E_P\!\left[b(X)\mathbb E_P\!\left[\mathbf 1\{A=a\}Y_a\mid X\right]\right]\\
&=\mathbb E_P\!\left[b(X)\mathbb E_P\!\left[\mathbf 1\{A=a\}\mid X\right]
                         \mathbb E_P\!\left[Y_a\mid X\right]\right]\\
&=\mathbb E_P\!\left[b(X)\rho_a(X)g_a(X)\right].
\end{aligned}
\]
Here the second equality is exactly the conditional factorization supplied by \(Y_a\perp A\mid X\). On the other hand, the defining property of the observed-arm regression \(\mu_a\), followed by conditioning on \(X\), gives
\[
\begin{aligned}
\mathbb E_P\!\left[b(X)\mathbf 1\{A=a\}Y\right]
&=\mathbb E_P\!\left[b(X)\mathbf 1\{A=a\}\mu_a(X)\right]\\
&=\mathbb E_P\!\left[b(X)\rho_a(X)\mu_a(X)\right].
\end{aligned}
\]
Consequently,
\[
\mathbb E_P\!\left[b(X)\rho_a(X)\{\mu_a(X)-g_a(X)\}\right]=0
\]
for every such \(b\). Taking \(b\) successively to be the indicators of the sets on which the \(X\)-measurable integrand is positive and negative shows that
\[
\rho_a(X)\{\mu_a(X)-g_a(X)\}=0
\qquad P\text{-almost surely}.
\]
Since \(\rho_a(X)\ge e_0>0\), division by \(\rho_a(X)\) is valid, and therefore
\[
\mu_a(X)=g_a(X)=\mathbb E_P[Y_a\mid X]
\qquad P\text{-almost surely}.
\]
Thus \(\mu_a\) is a version of \(x\mapsto\mathbb E_P[Y_a\mid X=x]\) for each \(a\in\{0,1\}\).

By linearity of conditional expectation and the two armwise identities just proved,
\[
\mathbb E_P[Y_1-Y_0\mid X]
=g_1(X)-g_0(X)
=\mu_1(X)-\mu_0(X)
\qquad P\text{-almost surely}.
\]
The representative convention defining \(\tau\) says that \(\tau=\mu_1-\mu_0\) \(P_X\)-almost everywhere, while \(\mathrm{ass{:}effect\mbox{-}holder}\) gives \(\tau\in\mathcal H^\gamma(L)\). Since \(\gamma>0\), this representative is continuous. It follows that \(\tau\) is a continuous representative of \(x\mapsto\mathbb E_P[Y_1-Y_0\mid X=x]\).

It remains to verify that the pointwise representative is unique. Let \(f\) and \(g\) be two continuous functions on \([0,1]^d\) that agree \(P_X\)-almost everywhere. By \(\mathrm{ass{:}uniform\mbox{-}design}\), they agree Lebesgue-almost everywhere. If \(f(x_*)\ne g(x_*)\) at some \(x_*\in[0,1]^d\), continuity of \(f-g\) supplies an \(r>0\) such that
\[
|f(x)-g(x)|>\frac12|f(x_*)-g(x_*)|>0
\quad\text{for every }x\in B(x_*,r)\cap[0,1]^d.
\]
The relative ball \(B(x_*,r)\cap[0,1]^d\) has strictly positive Lebesgue measure, including when \(x_*\) lies on the boundary of the cube. This contradicts equality almost everywhere. Hence \(f=g\) everywhere, proving uniqueness.

Finally, the observed law determines each \(\mu_a\) up to \(P_X\)-null sets and therefore determines the almost-everywhere contrast \(\mu_1-\mu_0\). The preceding uniqueness argument turns that almost-everywhere class into exactly one continuous function. In particular, any two causal laws satisfying the stated assumptions and inducing the same observed law have the same continuous contrast at \(x_0\). Therefore
\[
\theta(P)=\tau_P(x_0)
\]
is identified as the value at \(x_0\) of the unique continuous representative of \(x\mapsto\mathbb E_P[Y_1-Y_0\mid X=x]\), as claimed.
